\section{A dual space with coded norming functionals}
\label{sec:norming}

We begin the real counterexample by constructing a dual space $E$ with
coordinate predual $X_0$. Its infinite-dimensional subspaces will admit
many separated functionals, and those functionals will be combined into
an operator that changes the predual. The two tasks require different
features of the norm: averages along coded paths give separation, while
estimates for combinations of different paths make the operator bounded.
We specify these goals below before choosing the averaging parameters.

Everything from this section through the construction of the counterexample
is over the real field.  Write $\Gamma=\omega_1$ and fix a partition
\[
 \Gamma=\bigcup_{k=0}^{\infty} C_k
\]
into pairwise disjoint unbounded sets.  For $\alpha\in\Gamma$, let
$\operatorname{col}(\alpha)$ denote its color.  A finite nonempty set has
both an ordinal range and a color range: these are the intervals between
its least and greatest ordinal, and between its least and greatest color.
Finite supports are \emph{successive} if their ordinal ranges are strictly
successive.  They are \emph{mixed-successive} if their ordinal ranges and
their color ranges are both strictly successive.

An ordinal interval means any convex subset of $\Gamma$, including the
whole set.  A color interval means a convex subset of
$\{0,1,2,\ldots\}$.  A \emph{rectangle} is the intersection of an ordinal
interval with the set of coordinates whose colors lie in a color
interval.  These are the projection sets used below.

\subsection{Parameters and finite coordinate closures}

There is one pure family for every $k\ge0$, and one mixed family, denoted
by $*$.  For the pure family $k$, set
\[
 s_k=k+2,\qquad r_k=s_k',\qquad s_{k+1}'\le q_k<s_k',
\]
where $a'=a/(a-1)$ is the conjugate exponent.  For the mixed family set
\[
 s_*=2,\qquad r_*=2,\qquad s_1'\le q_*<2.
\]
Parameters belonging to different families are chosen independently.
Whenever a family is fixed, we omit its subscript and write $s,r,q$.
In particular,
\[
 1<q<r=s',\qquad q'>r,\qquad \frac1{q'}<\frac1s.
\]

Our aim is a norm on $c_{00}(\Gamma)$ whose completion $E$ has the
following properties. Write $e_\alpha$ for its coordinate vectors and
$e_\alpha^*$ for the coordinate functionals, and set
$X_0=[e_\alpha^*: \alpha<\omega_1]\subset E^*$.
We shall prove $E=X_0^*$ under the canonical evaluation map. For each
family $t\in\{0,1,\ldots,*\}$ we shall also construct a set $\mathcal P_t$ of infinite paths, each
defining a functional $P\in B_{E^*}$.
The properties to be established are:
\begin{enumerate}
\item On each color, every finitely supported scalar family satisfies the
coordinate upper estimate
\[
 \Big\|\sum_{\alpha\in C_k}b_\alpha e_\alpha\Big\|
 \le \|(b_\alpha)\|_{\ell_{s_k}}.
\]
\item If $Q:E^*\to E^*/X_0$ is the quotient map, then finite
combinations of distinct members of one path family satisfy
\[
 \Big\|\sum_{P\in\mathcal P_t}a_P QP\Big\|
 \le \|(a_P)\|_{\ell_{q_t}}.
\]
\item Every infinite-dimensional closed $F\subset E$ has, in a single
family, an uncountable set of path restrictions $P|_F$ separated by a
positive distance in $F^*$.
\end{enumerate}
The coordinate estimate is proved in this section, the estimate modulo
$X_0$ follows from the path analysis in Section~\ref{sec:paths}, and
Section~\ref{sec:branches} proves the separation assertion using the
prescribed-weight pairs of Section~\ref{pair:section}.

Here is why the exponents connect adjacent colors. Under CH we can label
all paths in pure family $k$ by distinct coordinates of $C_{k+1}$.
The first estimate makes evaluation at those coordinates an
$\ell_{s_{k+1}'}$ family for every $y\in E^*$. Since
$s_{k+1}'\le q_k$, these coefficients can be inserted into the second
estimate to form a bounded combination of path cosets. In the adjoint
direction the conjugate inequality $q_k'\le s_{k+1}$ makes the
corresponding coordinate series converge in $E$. The mixed family uses
$C_1$ in the same way. Section~\ref{sec:predual} carries out this
construction with small summable coefficients, then uses the separated
restrictions to exclude separable quotients of the resulting predual.

The remaining parameters specify the norming operations that will give
these properties. A weight-$j$ average of at most $L_j$ successive norming functionals
will use coefficient $1/m_j$. If each functional has coefficient
$\ell_r$ norm at most one, their disjoint supports give the bound
$\theta_j=L_j^{1/r}/m_j$ for the average. The integers $d_j$ enter only the finite-depth
estimate for prescribed-weight pairs; their inequalities are recorded
here so that all parameter choices are made together.

\begin{lemma}[Choice of parameters]\label{ch:parameters}
For each family there are strictly increasing sequences of positive
integers $(L_j)$ and $(m_j)$, and positive integers $(d_j)$, such that,
with
\[
 \theta_j=\frac{L_j^{1/r}}{m_j},\qquad
 c_0=\sum_{j=1}^{\infty}\theta_j,\qquad
 B_j=2\sum_{h<j}L_h,
\]
one has
\begin{gather}
 c_0<\frac18,\qquad
 \sum_{j\ge J}\theta_j\le4\theta_J,\qquad
 \sum_{j=1}^{\infty}\frac1{m_j}<\frac18,
 \label{ch:parameter-smallness}\\
 \frac{\log(L_j/m_j)}{\log L_j}\longrightarrow\frac1s,
 \label{ch:parameter-growth}\\
 (1/8)^{d_j}\le\theta_j,\qquad
 (1+B_j)^{d_j}\bigl(1+B_j+L_j^{1/q'}\bigr)
 \le\frac{L_j}{m_j}.
 \label{ch:parameter-depth}
\end{gather}
In particular, $L_j/m_j\to\infty$.
\end{lemma}

\begin{proof}
Put $\eta_j=(2^{-10-j})^r$ and choose
$L_j=\lfloor\eta_jm_j^r\rfloor$.  By taking $m_j$ large, arrange
\[
 \eta_j/2\le L_j/m_j^r\le\eta_j.
\]
It follows that
\[
 2^{-1/r}2^{-10-j}\le\theta_j\le2^{-10-j}.
\]
These bounds imply the first two inequalities in
\eqref{ch:parameter-smallness}, since
$2^{1+1/r}<4$.  Choose $d_j$ before choosing $m_j$, so that
$(1/8)^{d_j}\le2^{-1/r}2^{-10-j}$.
At step $j$, the numbers $B_j$ and $d_j$ are already fixed.  Because
\[
 \frac{L_j}{m_j}=\theta_jL_j^{1/s}
 \quad\hbox{and}\quad \frac1{q'}<\frac1s,
\]
the second inequality in \eqref{ch:parameter-depth} holds for all
sufficiently large choices of $m_j$.  Increase $m_j$ further to make both
integer sequences strictly increasing, to ensure
$1/m_j<2^{-j-4}$, and to make $|\log\theta_j|/\log L_j<1/j$.
The last condition gives \eqref{ch:parameter-growth}; its positive limit
also gives $L_j/m_j\to\infty$.
\end{proof}

We next supply finite coordinate sets on which a rational functional can
be recorded by positions, without recording its ordinal labels.
Encoding supports in this way follows the method of
\cite[Sections~2.3--2.4]{ArgyrosLopezAbadTodorcevic2005}; we prove the
coherence properties needed here directly.

\begin{lemma}[Coherent finite closures]\label{ch:closures}
For every $\beta<\omega_1$ there are nested finite sets
$F_p(\beta)\subset[0,\beta]$, $p\ge1$, containing $\beta$, exhausting
$[0,\beta]$, and satisfying
\begin{equation}
 \alpha\in F_p(\beta)
 \quad\Longrightarrow\quad
 F_p(\alpha)=F_p(\beta)\cap[0,\alpha].
 \label{ch:coherence}
\end{equation}
Consequently, two such closures with the same index $p$ that contain a
common coordinate $\alpha$ have the same ordered initial segment through
$\alpha$.
\end{lemma}

\begin{proof}
Construct the sets by induction on $\beta$.  At zero use $\{0\}$, and at
a successor append $\beta$ to every closure of its predecessor.  Suppose
$\beta$ is a nonzero limit.  Enumerate its predecessors as $(a_i)$, and
choose an increasing cofinal sequence $(\alpha_i)_{i\ge0}$ below $\beta$
such that $\alpha_i\ge a_h$ for $h\le i$.
Let $p_0=1$.  Recursively choose strictly increasing $p_i>1$, $i\ge1$,
so large that $F_{p_i}(\alpha_i)$ contains $\alpha_{i-1}$ and
$a_0,\ldots,a_i$.  For $p_i\le p<p_{i+1}$ set
\[
 F_p(\beta)=F_p(\alpha_i)\cup\{\beta\}.
\]
At a switch, coherence at $\alpha_i$ gives
\[
 F_{p_i}(\alpha_{i-1})
 =F_{p_i}(\alpha_i)\cap[0,\alpha_{i-1}],
\]
so the sets remain nested.  The enumerated predecessors are eventually
included, and coherence follows directly from coherence at $\alpha_i$.
The last assertion follows by applying \eqref{ch:coherence} to the common
coordinate.
\end{proof}

For $\alpha\le\beta$ put
\[
 \rho(\alpha,\beta)=\min\{p:\alpha\in F_p(\beta)\}.
\]
For a finite nonempty set $A$, write
$\max\rho(A)=\max\{\rho(\alpha,\beta):\alpha,\beta\in A,
\alpha\le\beta\}$.  If $p\ge\max\rho(A)$, then
$A\subset F_p(\max A)$.

We will repeatedly use the following precise form of position alignment.
Suppose $F_p(\beta)$ and $F_p(\gamma)$ contain $\alpha$, and two arrays
on these closures have identical entries at corresponding order positions.
If the first array is supported strictly before $\alpha$, then the arrays
define exactly the same finitely supported function on $\Gamma$.
Indeed, the positions through $\alpha$ coincide by
Lemma~\ref{ch:closures}, and every later entry of the first array, hence
also of the second array, is zero.  Equality of position codes alone does
not assert equality of ordinal supports; the common coordinate is the
additional input in this observation.

\subsection{Operations, paths and the norm}

All norming functionals below are finitely supported rational arrays on
$\Gamma$.  A Type~I expression of weight $j$ in a fixed family has the
form
\begin{equation}
 \phi=\frac1{m_j}\sum_{h=1}^{d}g_h,
 \qquad 1\le d\le L_j,
 \label{ch:type-i}
\end{equation}
with nonzero children $g_h$ having successive supports in the pure case
and mixed-successive supports in the mixed case.  The membership
conditions on the children are specified shortly.

Following the special-sequence method of~\cite{GowersMaurey1993},
the next weight will record the preceding rational functionals.
A finite \emph{path} in one family is a list of nonzero Type~I pieces
$\phi_1,\ldots,\phi_l$, successive in that family's order, with declared
weights $j_1,\ldots,j_l$ and integer metadata $p_1<\cdots<p_l$ satisfying
\[
 p_i\ge j_i,\qquad
 p_i\ge\max\rho\left(\bigcup_{h\le i}\operatorname{supp}\phi_h\right).
\]
The code of the prefix through $i$ records all $j_h,p_h$, $h\le i$, and
the zero-filled arrays of all $\phi_h$, $h\le i$, at the order positions
of $F_{p_i}(\beta_i)$, where
$\beta_i=\max\bigcup_{h\le i}\operatorname{supp}\phi_h$.
There are countably many such codes.  In each family fix an injection
$\sigma$ from the set of possible codes to the positive integers such
that its value is greater than the final metadata entry.  Such an
injection is obtained by enumerating the codes and successively choosing
unused sufficiently large integers.  The path rule is
\begin{equation}
 j_{i+1}=\sigma(\text{code of the prefix through }i).
 \label{ch:path-rule}
\end{equation}
The first weight is arbitrary.  All weights along a path are strictly
increasing.  A path includes its pieces, declared weights and metadata;
it is not merely the sum of its pieces.

For a pure path, a crop means restriction of its sum to an ordinal
interval.  For a mixed path it means restriction to a rectangle.  A
weight is \emph{active} in a cropped path if its piece remains nonzero
after the crop.  Path rules always refer to the original uncropped
pieces. The distinction between an uncropped sequence and its surviving
indices also appears in~\cite[Section~3]{ArgyrosArvanitakisTolias2006};
here the surviving indices are the active weights.

For each $k$, start $K_k$ with $0$ and the coordinate functionals
$\pm e_\alpha^*$, $\alpha\in C_k$.  Enlarge it in finite stages by the
following pure operations:
\begin{enumerate}
 \item Type~I expressions \eqref{ch:type-i} whose children belong to
 earlier stages of $K_k$;
 \item Type~II expressions
 \begin{equation}
  \sum_{b=1}^{h}a_b\,P_{I_b}
    \left(\sum_{i=1}^{l_b}\phi_{b,i}\right),
  \qquad a_b\in\mathbb Q,\quad \sum_{b=1}^{h}|a_b|^q\le1,
  \label{ch:type-ii}
 \end{equation}
 where each list $(\phi_{b,i})_{i=1}^{l_b}$ is a finite path, the
 Type~I children of all its pieces belong to earlier stages, and the
 cropped paths have pairwise disjoint active weight sets.
\end{enumerate}
Here $P_A$ denotes restriction to coordinates in $A$, on both arrays and
vectors as appropriate.  Empty cropped terms are discarded.  A
functional may have several expressions or declared weights; all
arguments use a specified valid expression.

Now start $K$ with all members of all the $K_k$ as terminal functionals,
and apply the same finite-stage rules with the mixed parameters,
mixed succession and rectangle crops.  Children in these mixed rules
belong to earlier stages of $K$.

\begin{lemma}[Norm and projections]\label{ch:norm}
The formula
\[
 \|x\|=\sup_{f\in K}|f(x)|\qquad(x\in c_{00}(\Gamma))
\]
defines a norm.  In its completion $E$, the coordinate vectors have norm
one, and projections onto ordinal intervals and rectangles are
contractive.  Every vector has countable support and is a norm limit of
finitely supported vectors.  The corresponding pure norm on $C_k$ is
given by $K_k$ and agrees with the full norm.
\end{lemma}

\begin{proof}
First note that each $K_k$ is symmetric and stable under ordinal crops,
and $K$ is symmetric and stable under rectangle crops.  For Type~I,
distribute the crop to the children and discard zeros; the remaining
children still have the required succession.  For Type~II, intersect
the existing crops with the new crop.  This changes neither the paths
nor their codes and only removes active weights.  These observations
prove stability by induction on the construction stage.  A rectangle
crop of a pure terminal either gives zero or leaves an ordinal crop
within the same color.

Every coefficient of every member of $K_k$ or $K$ has absolute value
at most one.  At a Type~I node, disjointness of children gives the
stronger bound $1/m_j$.  At a Type~II node, a given coordinate meets at
most one piece of each path.  Contributions from different paths have
distinct active weights, so their absolute sum is at most
$\sum_j1/m_j<1$.  Induction first establishes these bounds in the pure
families and then in the mixed family.  Consequently
\[
 \|x\|_\infty\le\|x\|\le\|x\|_1
 \qquad(x\in c_{00}(\Gamma)).
\]
Crop stability proves contractivity of the asserted projections.  Each
coordinate functional therefore extends continuously to $E$.  A
countable sequence of finite approximations shows that every vector
vanishes outside a countable set.  Moreover, $K$ remains norming on the
completion, by finite approximation.  Since its members use only
finitely many coordinates, those coordinates determine every vector.

It remains to compare the pure and full norms.  We claim that
\begin{equation}
 \{f|C_k:f\in K\}\subset\operatorname{absconv}(K_k),
 \label{ch:color-restriction}
\end{equation}
where zero is allowed in the absolutely convex hull.  It holds for
terminals.  At a mixed Type~I node at most one child meets $C_k$.
At a mixed Type~II node at most one piece of each path meets $C_k$,
and at most one child of that piece does so.  By induction the
restriction is an absolutely convex combination of pure functionals
whose total coefficient mass is at most
$\sum_j1/m_j<1$.  This proves \eqref{ch:color-restriction}.  The opposite
norm inequality follows from $K_k\subset K$.
\end{proof}

\begin{lemma}[Coordinate upper estimates]
\label{ch:color-bound}\label{ch:coordinates}
For each $k$ and every finitely supported scalar family
$(b_\alpha)_{\alpha\in C_k}$,
\begin{equation}
 \left\|\sum_{\alpha\in C_k}b_\alpha e_\alpha\right\|
 \le \|(b_\alpha)\|_{\ell_{s_k}}.
 \label{ch:ell-s}
\end{equation}
Thus the same coordinate series converges in $E$ whenever
$(b_\alpha)\in\ell_{s_k}(C_k)$.
\end{lemma}

\begin{proof}
Fix a pure family and show by induction that each member of $K_k$ has
coefficient $\ell_r$ norm at most one.  This holds at coordinates.  A
Type~I piece of weight $j$ has coefficient norm at most
$L_j^{1/r}/m_j=\theta_j$, by the disjoint supports of its children.
A cropped path with active weight set $A_b$ therefore has coefficient
norm at most
\[
 t_b=\left(\sum_{j\in A_b}\theta_j^r\right)^{1/r}.
\]
For a Type~II expression, the triangle inequality and H\"older's
inequality give
\[
 \left\|\sum_b a_b P_{I_b}\sum_i\phi_{b,i}\right\|_{\ell_r}
 \le \|(a_b)\|_{\ell_q}\|(t_b)\|_{\ell_{q'}}
 \le\|(t_b)\|_{\ell_r}
 \le\left(\sum_j\theta_j^r\right)^{1/r}<1.
\]
We used $q'>r$ and disjointness of the active weight sets.
Now apply H\"older's inequality to vectors on $C_k$ and use
Lemma~\ref{ch:norm}.  Applying \eqref{ch:ell-s} to finite tails proves
convergence of every asserted infinite series.
\end{proof}

\subsection{The coordinate predual}

The construction so far provides a Banach space $E$.  Its coordinate
functionals span a predual; proving this now will avoid any appeal to
bounded completeness later.

\begin{proposition}\label{ch:predual}
Let
\[
 X_0=\overline{\operatorname{span}}\{e_\alpha^*:\alpha<\omega_1\}
 \subset E^*.
\]
The canonical evaluation map $E\longrightarrow X_0^*$ is an isometric
isomorphism.
\end{proposition}

\begin{proof}
Each $f\in K$ is a finite linear combination of coordinate functionals
and has norm at most one on $E$.  Thus $K\subset B_{X_0}$ and the
evaluation map is isometric.

For surjectivity fix $g\in X_0^*$ and put
$a_\alpha=g(e_\alpha^*)$.  We construct, by transfinite induction on
$\gamma\le\omega_1$, a vector $x_\gamma\in E$ whose coordinates are
$a_\alpha$ for $\alpha<\gamma$ and zero otherwise.  Each constructed
prefix satisfies
\begin{equation}
 \|x_\gamma\|\le\|g\|.
 \label{ch:prefix-bound}
\end{equation}
Indeed, for finite $f\in K$,
\[
 f(x_\gamma)=g(P_{<\gamma}f),\qquad
 \|P_{<\gamma}f\|_{X_0}\le1,
\]
using crop stability.  This proves the bound whenever the prefix has
been constructed.

Start with $x_0=0$, and at a successor append $a_\gamma e_\gamma$.
Suppose $\lambda$ is a limit ordinal and all earlier prefixes exist.
If $(x_\gamma)_{\gamma<\lambda}$ were not norm Cauchy, there would be
$\delta>0$ and successive intervals $[\alpha_i,\beta_i)$ below
$\lambda$ for which
\[
 v_i=x_{\beta_i}-x_{\alpha_i},\qquad \|v_i\|\ge\delta.
\]
The intervals can be chosen recursively from the failure of the
Cauchy condition.  Consider first the case that for some color $k$
and some $\delta'>0$, infinitely many $v_i$ satisfy
$\|P_{C_k}v_i\|\ge\delta'$.  For any $j$, select $L_j$ of them.
By the pure/full norm equality and symmetry, choose pure normers
supported in their respective intervals, with evaluation at least
$\delta'/2$.  Their weight-$j$ Type~I average belongs to $K_k$.
Evaluating on a prefix that covers all the selected intervals gives
\[
 \|g\|\ge\frac{\delta' L_j}{2m_j},
\]
contradicting $L_j/m_j\to\infty$.

In the remaining case, $\|P_{C_k}v_i\|\to0$ for every fixed $k$.
Pass to a subsequence and choose successive finite color intervals
$D_i$ so that
\[
 \|P_{\operatorname{col}^{-1}(D_i)}v_i\|\ge\delta/2.
\]
Here is the recursive choice.  After the preceding upper color cutoff
has been fixed, its finitely many color projections tend to zero, so
choose the next $v_i$ with their sum of norms less than $\delta/4$.
For this individual vector the projection onto sufficiently high colors
has norm less than $\delta/4$, by finite approximation and
contractivity.  The intervening color interval has the desired norm.
Take rectangle-cropped normers of evaluations at least $\delta/4$.
Their supports are mixed-successive.  Selecting $L_j$ of them and
using a mixed Type~I average again contradicts
\eqref{ch:prefix-bound}, now with lower bound $\delta L_j/(4m_j)$.

Thus the prefixes are Cauchy at every limit, including $\omega_1$.
Completeness supplies their limit, with the required coordinates and
earlier projections, and \eqref{ch:prefix-bound} is preserved.  The
terminal vector $x_{\omega_1}$ agrees with $g$ on every coordinate
functional, hence on $X_0$ by density.  This proves surjectivity.
\end{proof}
